% GENERATED FILE. Edit canonical lessons, then run npm run generate:books.
% canonical-source-hash: fnv1a64:764f7286a8df6a05
% canonical-lessons: LA-C403-doctus, LA-C403-indoctus, LA-C403-eruditus, LA-C403-sciens, LA-C403-imperitus

\chapter{Learned and Unlearned}
\label{ch:la-learned-and-unlearned}
\hlchaptermodality{403}

\begin{chapteropening}
\textbf{By the end of this chapter:} \emph{I can call someone learned or unlearned, educated or knowing, and say that someone lacks experience.}

Complete the last lesson of chapter 403: I can call someone learned or unlearned, educated or knowing, and say that someone lacks experience.
\end{chapteropening}

\section[doctus, docta, doctum]{doctus, docta, doctum — learned, skilled}

\label{lesson:LA-C403-doctus}



\begin{quote}
Before the new one: say the Latin for to think over, then the Latin for to call to mind.
\end{quote}

\subsection*{You'll want to know: doctus, docta, doctum}
\textbf{doctus, docta, doctum} — \textquotedblleft{}learned, skilled\textquotedblright{}.

\begin{quote}
Vir doctus et prūdēns est.
\end{quote}

\begin{quote}
\emph{He is a learned and wise man.}
\end{quote}

\emph{doctus, docta, doctum}, the participle of \emph{doceō}: taught.

\begin{cousinweb}
\emph{doctus} is the participle of \emph{doceō} \textquotedblleft{}to teach\textquotedblright{} (Lewis \& Short): one who has been taught.
\end{cousinweb}

\begin{grammarlens}[title={What you've built}]
One more word for telling, explaining and describing at B1.
\end{grammarlens}

\subsection*{Your turn}
\begin{itemize}
\raggedright
  \item \emph{Say it:} \emph{doctus, docta, doctum}
  \item \emph{Say it:} \emph{doctus, docta, doctum}, once more
  \item \emph{Say it:} a sentence of your own with \emph{doctus, docta, doctum}
  \item \emph{Recall:} say the Latin for to think over, then the Latin for to call to mind, then say \emph{doctus, docta, doctum} again
\end{itemize}

\begin{tcolorbox}[breakable,colback=teal!4,colframe=teal!35!black,title={Before you move on}]
What does doctus, docta, doctum mean? (\textquotedblleft{}Learned, skilled\textquotedblright{}.) Say it once more.
\end{tcolorbox}
\section[indoctus, indocta, indoctum]{indoctus, indocta, indoctum — unlearned, untaught}

\label{lesson:LA-C403-indoctus}



\begin{quote}
Before the new one: say the Latin for to call to mind, then the Latin for learned.
\end{quote}

\subsection*{You'll want to know: indoctus, indocta, indoctum}
\textbf{indoctus, indocta, indoctum} — \textquotedblleft{}unlearned, untaught\textquotedblright{}.

\begin{quote}
Indoctus est, sed nōn stultus.
\end{quote}

\begin{quote}
\emph{He is unlearned, but not foolish.}
\end{quote}

\emph{indoctus, indocta, indoctum}: \emph{in-} not + \emph{doctus}.

\begin{cousinweb}
\emph{indoctus} is built from \emph{in-} \textquotedblleft{}not\textquotedblright{} and \emph{doctus} \textquotedblleft{}taught\textquotedblright{} (Lewis \& Short), the participle of \emph{doceō}.
\end{cousinweb}

\begin{grammarlens}[title={What you've built}]
One more word for telling, explaining and describing at B1.
\end{grammarlens}

\subsection*{Your turn}
\begin{itemize}
\raggedright
  \item \emph{Say it:} \emph{indoctus, indocta, indoctum}
  \item \emph{Say it:} \emph{indoctus, indocta, indoctum}, once more
  \item \emph{Say it:} a sentence of your own with \emph{indoctus, indocta, indoctum}
  \item \emph{Recall:} say the Latin for to call to mind, then the Latin for learned, then say \emph{indoctus, indocta, indoctum} again
\end{itemize}

\begin{tcolorbox}[breakable,colback=teal!4,colframe=teal!35!black,title={Before you move on}]
What does indoctus, indocta, indoctum mean? (\textquotedblleft{}Unlearned, untaught\textquotedblright{}.) Say it once more.
\end{tcolorbox}
\section[ērudītus, ērudīta, ērudītum]{ērudītus, ērudīta, ērudītum — educated, learned}

\label{lesson:LA-C403-eruditus}



\begin{quote}
Before the new one: say the Latin for learned, then the Latin for unlearned.
\end{quote}

\subsection*{You'll want to know: ērudītus, ērudīta, ērudītum}
\textbf{ērudītus, ērudīta, ērudītum} — \textquotedblleft{}educated, learned\textquotedblright{}.

\begin{quote}
Fēmina ērudīta multōs librōs lēgerat.
\end{quote}

\begin{quote}
\emph{The educated woman had read many books.}
\end{quote}

\emph{ērudītus, ērudīta, ērudītum}, the participle of \emph{ērudiō}. English \emph{erudite}.

\begin{cousinweb}
\emph{ērudītus} is the participle of \emph{ērudiō} \textquotedblleft{}to educate\textquotedblright{} (Lewis \& Short), from \emph{ē} and \emph{rudis} \textquotedblleft{}rough\textquotedblright{}. English \emph{erudite} comes from it (EtymDB).
\end{cousinweb}

\begin{grammarlens}[title={What you've built}]
One more word for telling, explaining and describing at B1.
\end{grammarlens}

\subsection*{Your turn}
\begin{itemize}
\raggedright
  \item \emph{Say it:} \emph{ērudītus, ērudīta, ērudītum}
  \item \emph{Say it:} \emph{ērudītus, ērudīta, ērudītum}, once more
  \item \emph{Say it:} a sentence of your own with \emph{ērudītus, ērudīta, ērudītum}
  \item \emph{Recall:} say the Latin for learned, then the Latin for unlearned, then say \emph{ērudītus, ērudīta, ērudītum} again
\end{itemize}

\begin{tcolorbox}[breakable,colback=teal!4,colframe=teal!35!black,title={Before you move on}]
What does ērudītus, ērudīta, ērudītum mean? (\textquotedblleft{}Educated, learned\textquotedblright{}.) Say it once more.
\end{tcolorbox}
\section[sciēns, scientis]{sciēns, scientis — knowing; on purpose; skilled (in)}

\label{lesson:LA-C403-sciens}



\begin{quote}
Before the new one: say the Latin for unlearned, then the Latin for educated.
\end{quote}

\subsection*{You'll want to know: sciēns, scientis}
\textbf{sciēns, scientis} — \textquotedblleft{}knowing; on purpose; skilled (in)\textquotedblright{}.

\begin{quote}
Sciēns hoc fēcit, nōn imprūdēns.
\end{quote}

\begin{quote}
\emph{He did this knowingly, not by mistake.}
\end{quote}

\emph{sciēns, scientis}, the present participle of \emph{sciō}: knowing, on purpose. With the genitive, skilled in: \emph{bellī sciēns}.

\begin{cousinweb}
\emph{sciēns} is the present participle of \emph{sciō} \textquotedblleft{}to know\textquotedblright{} (EtymDB; Lewis \& Short). English \emph{science} and \emph{conscience} come from its family.
\end{cousinweb}

\begin{grammarlens}[title={What you've built}]
One more word for telling, explaining and describing at B1.
\end{grammarlens}

\subsection*{Your turn}
\begin{itemize}
\raggedright
  \item \emph{Say it:} \emph{sciēns, scientis}
  \item \emph{Say it:} \emph{sciēns, scientis}, once more
  \item \emph{Say it:} a sentence of your own with \emph{sciēns, scientis}
  \item \emph{Recall:} say the Latin for unlearned, then the Latin for educated, then say \emph{sciēns, scientis} again
\end{itemize}

\begin{tcolorbox}[breakable,colback=teal!4,colframe=teal!35!black,title={Before you move on}]
What does sciēns, scientis mean? (\textquotedblleft{}Knowing; on purpose; skilled (in)\textquotedblright{}.) Say it once more.
\end{tcolorbox}
\section[imperītus, imperīta, imperītum]{imperītus, imperīta, imperītum — inexperienced, unskilled}

\label{lesson:LA-C403-imperitus}



\begin{quote}
Before the new one: say the Latin for educated, then the Latin for knowing.
\end{quote}

\subsection*{You'll want to know: imperītus, imperīta, imperītum}
\textbf{imperītus, imperīta, imperītum} — \textquotedblleft{}inexperienced, unskilled\textquotedblright{}.

\begin{quote}
Imperītus nauta nāvem frēgit.
\end{quote}

\begin{quote}
\emph{The unskilled sailor wrecked the ship.}
\end{quote}

\emph{imperītus, imperīta, imperītum}: the opposite of the word for experienced, with \emph{in-}, not. With the genitive: \emph{bellī imperītus}, without experience of war.

\begin{cousinweb}
\emph{imperītus} is built from \emph{in-} \textquotedblleft{}not\textquotedblright{} and \emph{perītus} \textquotedblleft{}experienced\textquotedblright{} (EtymDB; Lewis \& Short), which comes from the root of \emph{experior} \textquotedblleft{}to try, to test\textquotedblright{}.
\end{cousinweb}

\begin{grammarlens}[title={What you've built}]
One more word for telling, explaining and describing at B1.
\end{grammarlens}

\subsection*{Your turn}
\begin{itemize}
\raggedright
  \item \emph{Say it:} \emph{imperītus, imperīta, imperītum}
  \item \emph{Say it:} \emph{imperītus, imperīta, imperītum}, once more
  \item \emph{Say it:} a sentence of your own with \emph{imperītus, imperīta, imperītum}
  \item \emph{Recall:} say the Latin for educated, then the Latin for knowing, then say \emph{imperītus, imperīta, imperītum} again
\end{itemize}

\begin{tcolorbox}[breakable,colback=teal!4,colframe=teal!35!black,title={Before you move on}]
What does imperītus, imperīta, imperītum mean? (\textquotedblleft{}Inexperienced, unskilled\textquotedblright{}.) Say it once more.
\end{tcolorbox}
